% Intentionally flawed example manuscript for testing the deterministic checks.
% Every flaw below is annotated with the rule it violates.
\documentclass{article}
\begin{document}

\title{A Novel Framework For Advanced Simulation Of Materials}  % LG-4: Title Case
\maketitle

\begin{abstract}
In order to simulate the the mechanical behavior of magnesium alloys under
complex loading conditions, we present a comprehensive computational framework
based on a state based peridynamic formulation combined with machine learning
surrogates trained on high fidelity finite element simulations using adaptive
mesh refinement with residual based error estimators, where the coupling is
realized through a partitioned scheme with dynamic relaxation, and we further
investigate the influence of the horizon size, the mesh density, the time step
size, the boundary conditions, the loading rate, the temperature, and several
other parameters on the accuracy and stability of the proposed approach, and
moreover we discuss implementation aspects in detail including data structures,
parallelization strategies, load balancing, and memory layout considerations,
and finally we also report preliminary results for a number of additional test
cases which were not the main focus of this work but which we include for
completeness, demonstrating that the framework can in principle also be applied
to thermo-mechanical problems, fatigue, and corrosion, although a detailed
investigation of these topics is left for future work, since the present paper
focuses primarily on the quasi-static regime, and the results show good
agreement with reference solutions in most cases considered here, while we
additionally provide an extensive discussion of the calibration procedure for
the surrogate models, the selection of hyperparameters such as network depth,
width, activation functions, learning rate schedules, and regularization
strategies, the generation of training data including sampling strategies in
parameter space, and the treatment of boundary effects near the surfaces of
the computational domain, which require special correction procedures in
nonlocal models, and we conclude with an outlook on possible extensions of
the framework towards multiscale and multiphysics applications in the context
of digital twins of engineering structures.
\end{abstract}
% ST-1: far beyond 250 words of technical detail, no clean problem/solution arc

\section{Introduction}  % ST-2: no novelty paragraph at the end

The simulation of fracture is important. In order to understand fracture,
many methods have been developed. The results in Figure 1 show that classical
methods have difficulties with discontinuities.  % LG-3, LG-7, TX-3 (hardcoded ref)

It can be seen that peridynamics is a promising alternative. Previous studies
have shown its potential. The method was used by many authors and it shows a
nonlinear stress.  % LG-7 scaffolding; LG-1 "nonlinear stress" (semantic, LLM pass)

The paper is organized as follows. See Figure 2 for an overview.  % LG-7

\section{Methods For Peridynamic Simulation}  % LG-4 Title Case, ST-5 vs. next section

% TODO: rewrite this section after Arman's comments  % TX-1
The deformation is governed by the equation of motion, where rho denotes
density.

% Old derivation, kept just in case:
% The bond force density f depends on the stretch s as follows.
% f(s) = c s if s < s_c and zero otherwise, where c denotes the
% micromodulus constant which can be calibrated against the bulk
% modulus K of classical elasticity theory.
% This was replaced by the state based formulation below.
% TX-1: commented-out block

We use a one dimensional model with a stress strain relation that is rate
dependent.  % LG-5: missing hyphens (modifier use)

\subsection{discretization}  % LG-4: lowercase heading start (style break, ST-5)

Fig. 3 shows the discretization. The convergence is shown in Figure 4.
% TX-2: "Fig." vs "Figure" mixed; TX-3 hardcoded

\section{Results}

The error decreases quickly with mesh refinement. The agreement with the
the reference solution is excellent.  % FS-3 unquantified; LG-6 doubled "the"

\begin{figure}
  \centering
  % \includegraphics{convergence.pdf}
  \caption{Convergence of the proposed method for the one dimensional test
  case described in Section 3, where the error is measured in the L2 norm
  against the analytical solution, computed on a sequence of five uniformly
  refined discretizations, with the horizon held fixed in physical units,
  demonstrating second order accuracy as expected from the theory.}
  % FG-4: caption far too long; LG-5 "one dimensional", "second order"
  \label{fig:convergence-unused}  % FG-5: label never referenced
\end{figure}

\section{Conclusion}

In order to summarize, we presented a framework.  % LG-3

\end{document}
